\section{Overview}

\adnote[title={REVISED TIMELINE — April 22, 2026}]{This timeline has been revised from the original proposal to reflect an accelerated deposit target of October 2026, driven by a postdoctoral application requiring degree completion by that date. The committee chair/advisor has reviewed and confirmed this accelerated trajectory as feasible. This is sot to say this \textit{will} happen this way for certain, but it is the current plan in place.  

Key scope changes are documented in Section~\ref{sec:scope-changes}.}

The timeline presented here outlines the remaining trajectory of the dissertation research and writing from the present (April 2026) through ProQuest deposit (October 2026). Proposal/IRB approval was received in mid–late February 2026; given this the research period is now well underway.

The timeline is structured around two compressed phases that honor the cyclical logic of Speculocultural Technopoiesis (ST) while enforcing the disciplined forward momentum required.

\textbf{Revised Target Defense Date}: September 2026 (with deposit by October 2026)

\textbf{Original Target Defense Date}: December 2026 (revised upward)

\textbf{Remaining Research Period}: Approximately 6 months (April – October 2026)

\section{Scope Changes from Original Proposal}
\label{sec:scope-changes}

The following strategic scope reductions have been made to make the October 2026 deposit achievable without compromising the intellectual integrity and authenticity of the work. These changes have been discussed with my advisor/chair.

\subsection*{Primary Movement (Deep Analysis)}
\textbf{Movement II: Water — \textit{Did Sun Ra Fear Cryosleep?}} is designated the primary movement and the central site of deep theoretical analysis, collaborative practice documentation, and conference presentation. This movement has the richest existing literary and theoretical foundation and is the current making focus. It anchors the dissertation's demonstration of Speculocultural Technopoiesis.

\subsection*{Adjacent Movements (Preliminary / In Development)}
\textbf{Movement I: Earth — \textit{Buffalo Soldier}} and \textbf{Movement III: Skin — \textit{Don't Play with My Hair}} are documented as compositional works in development: architecturally present and theoretically grounded. They should be understood as a continuation of practice that extends beyond the dissertation. The dissertation argues for the framework through deep engagement with Movement II and preliminary embodiment in Movements I and III, framing the latter as an open score. Both intentionally unfinished, as research-creation methodology allows for.

\subsection*{Collaborative Workshop Scope}
\begin{itemize}[noitemsep]
	\item \textbf{Maximum participants}: 3 (at least 1 confirmed)
	\item \textbf{Maximum workshops}: 3 sessions total
	\item \textbf{IRB}: Already approved; this scope is within the approved protocol
\end{itemize}

\subsection*{Writing Approach}
Weekly writing blocks are non-negotiable from this point forward, integrated with making rather than deferred until after completion. Which is exactly how research-creation and critical technical practice suggests work to be done.  The 155-page proposal provides a solid translation scaffold (as it was originally intended). Individual chapter work proceeds by expansion and editorial transformation from the proposal case studies, instead of from scratch.

\section{Timeline Visualization}

\subsection*{Remaining Timeline at a Glance (April – October 2026)}
% GANTT CHART — Revised for October 2026 deposit
% Columns: Apr | May | Jun | Jul | Aug | Sep | Oct  (7 months)
\begin{ganttchart}[
	hgrid,
	vgrid,
	x unit=1.1cm,
	y unit title=0.6cm,
	y unit chart=0.55cm,
	title height=1
	]{1}{7}
	\gantttitle{2026}{7} \\
	\gantttitle{Apr}{1}
	\gantttitle{May}{1}
	\gantttitle{Jun}{1}
	\gantttitle{Jul}{1}
	\gantttitle{Aug}{1}
	\gantttitle{Sep}{1}
	\gantttitle{Oct}{1} \\

	% Phases
	\ganttbar[bar/.append style={fill=blue!20}]{Phase A: Deep Making}{1}{4} \\
	\ganttbar[bar/.append style={fill=green!20}]{Phase B: Synthesis \& Writing}{3}{7} \\

	% Making cycles (Movement II primary; I and III light)
	\ganttbar[bar/.append style={fill=purple!40}]{Mvmt II: Making Cycle 1}{1}{2} \\
	\ganttbar[bar/.append style={fill=purple!40}]{Mvmt II: Making Cycle 2}{3}{4} \\
	\ganttbar[bar/.append style={fill=purple!15}]{Mvmt I \& III: Light Dev.}{2}{4} \\

	% Collaborative workshops
	\ganttbar[bar/.append style={fill=orange!30}]{Workshop 1}{3}{3} \\
	\ganttbar[bar/.append style={fill=orange!30}]{Workshop 2}{4}{4} \\
	\ganttbar[bar/.append style={fill=orange!30}]{Workshop 3}{5}{5} \\

	% Writing (continuous, not deferred)
	\ganttbar[bar/.append style={fill=gray!30}]{Weekly Writing (all mvmts)}{1}{7} \\

	% Conferences
	\ganttmilestone[milestone/.append style={fill=red}]{SfNC + UCR Deadline}{1} \\
	\ganttmilestone[milestone/.append style={fill=red}]{SfNC Conference}{2} \\
	\ganttmilestone[milestone/.append style={fill=red}]{UCR Co-Creativity}{2} \\

	% Committee and deposit milestones
	\ganttmilestone{Committee Check-in I (May)}{1} \\
	\ganttmilestone{Committee Check-in II (May)}{3} \\
	\ganttmilestone{Committee Check-in III (May)}{5} \\
	\ganttmilestone{Full Draft to Committee}{5} \\
	\ganttmilestone{Defense}{6} \\
	\ganttmilestone[milestone/.append style={fill=black}]{ProQuest Deposit}{7}
\end{ganttchart}

\section{Milestones}

\adnote[title={Status as of April 2026}]{The phases below reflect the \textbf{remaining} work. Earlier work (proposal/IRB approval, technical infrastructure, initial making cycles, and the first collaborative relationship) is complete.}

\subsection*{Phase A: Deep Making and Conference Preparation (April – July 2026)}
\textit{4 months | Focus: Primary making cycles for Movement II; conference deliverables; workshops 1–3; concurrent weekly writing}

\textbf{April 2026 (current)}
\begin{itemize}[noitemsep]
	\item Advisory meeting: present revised scope and timeline (this document)
	\item Finalize Movement II Making Cycle 1: hardware/sensor integration for conference demo
	\item UCR Co-Creativity video: produce process/concept video for conference submission (deadline: May 1)
	\item Begin weekly writing blocks: Composer's Note and Movement II opening sections (editorial translation from proposal)
	\item \textit{Deliverable}: Advisory approval of revised scope; conference video draft
\end{itemize}

\textbf{May 2026}
\begin{itemize}[noitemsep]
	\item \textbf{Conference deadlines (May 1)}: SfNC submission + UCR Co-Creativity submission (video link, updated abstract, bio)
	\item Committee check-in: share revised timeline, Movement II progress, scope rationale
	\item Movement II Writing: \textit{Entry: The Question} and \textit{Log 1: The Wake} sections (full drafts from proposal scaffolding)
	\item Movements I \& III: light development = review existing placeholder sections, identify what requires new writing vs. editorial expansion
	\item \textit{Deliverable}: Submitted conference materials; Workshop 1 documentation; 2 full movement sections drafted
\end{itemize}

\textbf{June 2026}
\begin{itemize}[noitemsep]
	\item \textbf{SfNC Conference}: Philadelphia (June 2) - present Movement II work
	\item \textbf{UCR Co-Creativity Conference}: Riverside, CA (June 5–6) - present Movement II work
	\item Workshop 1 with confirmed collaborator(s): document via audio/video, field notes, reflective journal
	\item Movement II Writing: \textit{Log 2: What Water Holds} and \textit{Log 3: Temporal Escape} sections
	\item Begin Composer's Note chapter: full editorial draft from proposal's framework sections
	\item \textit{Deliverable}: Conference presentations complete; 4 movement sections drafted; Composer's Note draft
\end{itemize}

\textbf{July 2026}
\begin{itemize}[noitemsep]
	\item Movement II Making Cycle 2: refinement toward final installation configuration
	\item Workshop 2 (follow-up session): document fully
	\item Movement II Writing: \textit{Log 4: Instrument Built} and \textit{Dissolution} sections - complete Movement II
	\item Movement I Writing: Ground chapter - editorial expansion of existing Buffalo Soldier sections
	\item Movement III Writing: Touch chapter - editorial expansion of existing Don't Play with My Hair sections
	\item \textit{Deliverable}: Movement II writing complete; Movements I \& III in near-complete draft; all workshop documentation complete
\end{itemize}

\subsection*{Phase B: Synthesis, Writing Completion, and Defense Preparation (August – October 2026)}
\textit{3 months | Focus: Complete all remaining writing; committee review; defense; deposit}

\textbf{August 2026}
\begin{itemize}[noitemsep]
	\item Workshop 3 (final session): document fully; begin synthesis across workshop data into chapters
	\item Complete all remaining chapters: Coda, Note on Structure, Appendices, Preliminary matter (Abstract, Acknowledgements)
	\item Full dissertation assembly: integrate all movements, front matter, back matter
	\item Internal review: proofread, check cross-references, finalize figures and technical documentation
	\item \textbf{Submit complete draft to committee} by end of August
	\item \textit{Deliverable}: Complete draft submitted to committee for review
\end{itemize}

\textbf{September 2026}
\begin{itemize}[noitemsep]
	\item Committee reads and returns feedback (target: 3–4 week turnaround)
	\item Revisions based on committee feedback
	\item Prepare defense presentation and supporting materials
	\item \textbf{Dissertation defense} (target: mid–late September 2026)
	\item \textit{Deliverable}: Successful defense; final revision list from committee
\end{itemize}

\textbf{October 2026}
\begin{itemize}[noitemsep]
	\item Final revisions and formatting (ProQuest compliance check)
	\item Graduate school submission and approval
	\item \textbf{ProQuest deposit}: target end of October 2026
	\item \textit{Deliverable}: Deposited dissertation; degree conferred
\end{itemize}

\subsection*{Contingency Planning}

\begin{itemize}[noitemsep]
	\item \textbf{Writing pace}: If a writing week is lost (conference travel, illness), the next week carries a double target. Goal: no weeks off the page
	\item \textbf{Collaborative scheduling}: If workshop participants become unavailable, autoethnographic solo practice is methodologically sufficient per the approved IRB protocol; the dissertation does not depend on a minimum number of workshop participants beyond what is already confirmed
	\item \textbf{Making vs. writing tension}: If time pressure forces a choice, writing takes priority over further making iterations after July 2026; the dissertation's contribution is the framework demonstrated through the work, not the final installation's polish
\end{itemize}

\subsection*{Alignment with Institutional Requirements (Revised)}

\begin{itemize}[noitemsep]
	\item \textbf{Proposal/IRB approval}: Complete (mid–late February 2026)
	\item \textbf{Candidacy status}: Achieved
	\item \textbf{Committee check-in}: May 2026 (scope revision presentation)
	\item \textbf{Complete draft submission}: August 2026 (4–6 weeks before defense)
	\item \textbf{Final defense}: September 2026
	\item \textbf{ProQuest deposit}: October 2026
\end{itemize}

\subsection*{Timeline Summary Table}

\begin{table}[H]
	\centering
	\small
	\begin{tabular}{|l|p{3cm}|p{4.5cm}|p{4.5cm}|}
		\hline
		\textbf{Phase} & \textbf{Timeframe} & \textbf{Key Activities} & \textbf{Deliverables} \\
		\hline
		\rowcolor{blue!10}
		\textbf{Phase A} & Apr – Jul 2026 (4 months) &
		\begin{itemize}[nosep,leftmargin=*]
			\item Mvmt II making cycles 1–2
			\item 3 collaborative workshops
			\item Conference deliverables (SfNC, UCR)
			\item Weekly writing throughout
		\end{itemize} &
		\begin{itemize}[nosep,leftmargin=*]
			\item Conference presentations
			\item Mvmt II fully drafted
			\item Mvmts I \& III near-complete
			\item Most workshops documented
		\end{itemize} \\
		\hline
		\rowcolor{green!10}
		\textbf{Phase B} & Aug – Oct 2026 (3 months) &
		\begin{itemize}[nosep,leftmargin=*]
			\item Complete all chapters
			\item Committee review
			\item Defense preparation
			\item Final revisions
		\end{itemize} &
		\begin{itemize}[nosep,leftmargin=*]
			\item Full draft (August)
			\item Defense (September)
			\item ProQuest deposit (October)
		\end{itemize} \\
		\hline
		\rowcolor{gray!10}
		\textbf{Ongoing} & Apr – Oct &
		\begin{itemize}[nosep,leftmargin=*]
			\item Weekly writing blocks
			\item Field documentation
			\item Advisor communication
		\end{itemize} &
		Continuous prose output \\
		\hline
	\end{tabular}
	\caption{Revised phase-based timeline for October 2026 deposit.}
	\label{tab:timeline-revised}
\end{table}